% ======================================================================
% SECTION 3: RESULTS
% Four subsections, one per simulation. Each subsection has the same
% internal layout: (i) what was generated, (ii) advantages, (iii)
% disadvantages, (iv) one or more comprehensive ASCII diagrams to embed
% verbatim with \begin{verbatim} ... \end{verbatim}, and (v) hour-level
% or stage-level patient-and-robot examples to call out by name.
% ======================================================================

\subsection{Simulation 1: Continuous National 24/7 Real-Time Clinical Trial}
\label{sec:results-sim1}

[PROCESSING INSTRUCTION (Simulation 1, 5 to 8 paragraphs): Read the
README at new-trial/national-24-7-trial/README.md and every per-hour file
in new-trial/national-24-7-trial/hour-00/ through hour-55/ plus
new-trial/national-24-7-trial/extra-hours/hour-56/ through hour-83/
(seven files per hour: simulation.md, robot\_logs.md, patient\_records.md,
psl\_scores.md, diagram\_facility.txt, diagram\_patient\_flow.txt,
diagram\_robot\_status.txt). Read the final-commit summary at
new-trial/final-commit/final\_24h\_summary.md.

Write a results subsection that opens by stating the scope: 84 hours of
continuous trial coverage (hour-00 through hour-83) generated in response
to the FDA 28 April 2026 RTCT announcement, with 4 sites (Houston,
Philadelphia, Boston, Texas Medical Center) and 116 robot instances
streaming endpoints to a Paradigm Health-style aggregator, then to the
FDA real-time API. Quote the C-PSL (Continuity-PSL) rolling 24-hour
metric introduced for continuous trials.

Include the following sample comprehensive ASCII diagrams verbatim:
\begin{itemize}[leftmargin=*]
\item hour-00: hour\_00\_diagram\_facility.txt (network kickoff state)
\item hour-12: hour\_12\_diagram\_facility.txt (mid-day peak utilization)
\item hour-12: hour\_12\_diagram\_patient\_flow.txt (cross-site flow)
\item hour-23: hour\_23\_diagram\_robot\_status.txt (end-of-day-1 robots)
\item hour-47: hour\_47\_diagram\_facility.txt (day-2 close, comparison)
\item extra-hours/hour-72: hour\_72\_diagram\_facility.txt (day-3 close)
\item final-commit/final\_diagram\_facility.txt (cumulative summary)
\end{itemize}

Identify the following individual patient and robot examples by name and
hour, drawn from the per-hour patient\_records.md and robot\_logs.md files:
\begin{itemize}[leftmargin=*]
\item Hour 00, SITE-A patient PT-A-001 arrival and PSL initialization,
plus RTPOS-01/TRACK-01 robot pair startup status.
\item Hour 12, SITE-A four active sessions (RTPOS-03/TRACK-03 wrap,
RTPOS-01/TRACK-01, IMAGE-04) with AE-002 resolution at 12:30 and the
Paradigm Health aggregator state ("All UP. 8 endpoints + 1 AE-resolution").
\item Hour 23, end-of-day-1 wrap with the cumulative C-PSL rolling
metric and any unresolved safety signals.
\item Hour 47, day-2 close with cross-site comparison and any cumulative
PSL improvement vs hour-23.
\item extra-hours/hour-83, the last approximated-diagram hour, noting in
the prose that this hour is in the extra-hours subdirectory due to
diagram approximation.
\end{itemize}

Advantages to highlight (one paragraph): repository-scale processing -
Claude Code Opus 4.7 Max held all 168 patients across 4 sites and 116
robots in a single inference window, generated 3 ASCII diagrams plus 4
markdown files per hour without iterative human prompting, and emitted
hourly real-time commits to GitHub at a 1-hour cadence (24 commits per
day, 84 commits across the run). The minute-resolution per-hour
simulation.md file produced finer-grained patient and robot timelines
than any narrow supervised model can produce.

Disadvantages (one paragraph): no separate code agents for local compute
were provided in this simulation; the entire output is text-only (markdown
plus ASCII), so the work is computationally cheap but offers no Python
execution surface for downstream automation. Hours 56 through 83 were
relocated to extra-hours/ because the diagrams contain approximations
rather than minute-precise computations.

Cite \cite{kawchak_2026_19176370} and \cite{fda2026realtime}.]

\subsection{Simulation 2: Single-Patient 10-Stage Patient Journey}
\label{sec:results-sim2}

[PROCESSING INSTRUCTION (Simulation 2, 4 to 6 paragraphs): Read the
README at patient-journey/README.md, the prior paper at
patient-journey/paper/patient\_journey\_paper.tex, and every Python stage
script: patient-journey/stage\_01\_prescreening.py through
stage\_10\_closeout.py, plus the orchestrators
patient-journey/master\_journey.py and patient-journey/patient\_state.py.
Read the deliverables in patient-journey/deliverables/ and the diagrams
in patient-journey/diagrams/.

Write a results subsection that summarizes the journey of patient
PAT-2026-0042, a 58-year-old female with Stage IIIB NSCLC adenocarcinoma
(ECOG 1, PD-L1 65 percent, TMB 14 mut/Mb), through the 10 sequential
stages: pre-screening, enrollment, digital twin initialization, robot
qualification, robotic surgery, post-operative recovery, immunotherapy,
federated learning, surveillance, and trial closeout. Reference the FDA
cost-savings analysis (30 to 50 percent trial cost reductions, \$390M to
\$650M against a \$1.3B baseline) and the regulatory adaptations (21 CFR
312, 21 CFR 50, ICH E6(R3)).

Include the following sample comprehensive ASCII diagrams verbatim:
\begin{itemize}[leftmargin=*]
\item patient-journey/diagrams/timeline\_progress.txt (timeline view)
\item patient-journey/diagrams/regulatory\_progress.txt (regulatory view)
\item patient-journey/diagrams/clinical\_progress.txt (clinical view)
\item patient-journey/deliverables/ FDA cost-savings diagram
\end{itemize}

Identify the following individual stage-level events by name:
\begin{itemize}[leftmargin=*]
\item Stage 5 (stage\_05\_surgery.py): da Vinci Xi robotic lobectomy at
USL 87.5, 168-minute procedure, R0 resection with negative margins.
\item Stage 7 (stage\_07\_immunotherapy.py): 35 pembrolizumab cycles
managed by Franka Emika at USL 88.75, complete response with recurrence
risk declining from 35 percent to 3 percent over 36 months of event-free
survival.
\item Stage 9 (stage\_09\_surveillance.py): federated learning surveillance
state and any per-cycle digital twin updates.
\end{itemize}

Advantages (one paragraph): single-patient depth - this simulation
produces the most regulatory-grounded patient prediction in the four,
because every stage is mapped to a specific 21 CFR / ICH section, the
cost-savings analysis ties safety and efficacy predictions to dollar
figures, and the journey produces 12 Python modules accessible at every
commit. Real-time human monitorability is documented by timestamped text
diagrams.

Disadvantages (one paragraph): the simulation is restricted to a single
patient and is not as extensive as the multi-patient continuous trial in
Simulation 1. There is no national network model and no multi-site
aggregation. The single-patient orientation limits population-level
prediction calibration to TRIPOD+AI / CREMLS standards.

Cite \cite{kawchak_2026_19119939}, \cite{Collins2024TRIPODAI}, and
\cite{ElEmam2024CREMLS}.]

\subsection{Simulation 3: 24-Hour Autonomous Sponsor with Many Agents}
\label{sec:results-sim3}

[PROCESSING INSTRUCTION (Simulation 3, 5 to 7 paragraphs): Read the
README at sponsor/final\_paper/README.md and the script directories:
sponsor/final\_paper/scripts/core\_agents/ (53 agent .py files plus
\_\_init\_\_.py), sponsor/final\_paper/scripts/coordination/ (4
coordination modules), sponsor/final\_paper/scripts/dashboard/,
sponsor/final\_paper/scripts/diagrams/ (75 ASCII diagrams),
sponsor/final\_paper/scripts/hourly/sponsor\_hour\_00.py through
sponsor\_hour\_23.py, sponsor/final\_paper/scripts/safety/, and
sponsor/final\_paper/scripts/sponsor\_server/. Read the JSON outputs at
sponsor/final\_paper/scripts/hourly/output/sponsor\_hour\_00.json through
sponsor\_hour\_23.json and the cumulative
sponsor/final\_paper/scripts/output/sponsor\_24h\_summary.json.

Write a results subsection that summarizes the 24-hour autonomous sponsor
operating system. State the agent count (53 core agents organized into 4
layers: governance, study execution, site/robotics, trust), the script
count (24 hourly scripts, 75 ASCII diagrams, 24 JSON outputs), and the
PSL trajectory across the day.

Include the following sample comprehensive ASCII diagrams verbatim:
\begin{itemize}[leftmargin=*]
\item sponsor/final\_paper/scripts/diagrams/agent\_workload\_hour\_00.txt
(start of day workload by agent)
\item sponsor/final\_paper/scripts/diagrams/agent\_workload\_hour\_12.txt
(mid-day peak; this file is 34 tasks total across 12 agents, with
clinops\_agent at 5 tasks and clinical\_accountability\_agent at 4 tasks
and 1 escalation)
\item sponsor/final\_paper/scripts/diagrams/agent\_workload\_hour\_23.txt
(end of day workload)
\item one cumulative diagram drawn from sponsor\_24h\_summary.json
\end{itemize}

Identify the following individual hour-level patient and robot examples
by name:
\begin{itemize}[leftmargin=*]
\item Hour 00, overnight, 2 patients: name the two patients from
sponsor\_hour\_00.json and the robot authorization decisions made by
robot\_execution\_gateway.
\item Hour 12, mid-day peak, 5 patients: identify the safety\_agent
escalation and the clinical\_accountability\_agent decision flow.
\item Hour 23, overnight, 2 patients: identify the day-end audit trail
manager final state.
\end{itemize}

Advantages (one paragraph): the sponsor simulation includes the largest
number of agents in the four-simulation set (53 core agents), the largest
number of ASCII diagrams (75), and an end-to-end FastAPI sponsor server
under sponsor/final\_paper/scripts/sponsor\_server/ that exposes 5 API
endpoints. The 24 Python scripts plus 24 JSON files give downstream
analytics a structured execution surface that Simulation 1 lacks.

Disadvantages (one paragraph): the simulation lacks some of the advanced
detail of the continuous-trial Simulation 1, notably the minute-resolution
per-hour simulation.md narrative. The 24-hour scope, while complete, does
not exercise multi-day continuity that Simulation 4 introduces.

Cite \cite{kawchak_2026_19396256}.]

\subsection{Simulation 4: 168-Hour 7-Day Extension with Local Verification}
\label{sec:results-sim4}

[PROCESSING INSTRUCTION (Simulation 4, 6 to 9 paragraphs): Read the
README at sponsor/final\_paper/168\_hours/README.md and the day-level
directories sponsor/final\_paper/168\_hours/day\_01/ through
sponsor/final\_paper/168\_hours/day\_07/. Each day directory contains
hourly/sponsor\_hour\_NNN.py (24 scripts per day, 168 total),
hourly/output/sponsor\_hour\_NNN.json (168 JSON outputs total),
diagrams/agent\_workload\_hour\_NNN.txt (75 per day, 525 total), and a
day-level summary. Read the local verification README at
sponsor/final\_paper/168\_hours/instructions/core\_i5\_6200u\_4gb/README.md
and the verification output at
sponsor/final\_paper/168\_hours/instructions/core\_i5\_6200u\_4gb/sponsor\_168h\_summary.json.
Also read the orchestrators _commit_hour.sh, \_config.py, \_gen\_day\_summary.py,
\_gen\_hourly.py, \_gen\_init.py, and run\_168h\_simulation.py at
sponsor/final\_paper/168\_hours/.

Write a results subsection that summarizes the 7-day extension. State
the cumulative metrics from the 168-hour summary: 168 hours, 2,016
sponsor decisions, 1,336 patients processed, 125 escalations, 1,336 robot
authorizations, PSL start 63.4 to PSL end 70.0 (improvement +6.6), 525
text diagrams, 168 Python scripts, 168 JSON output files, 7 daily summary
files, 7 branches, 168 GitHub commits, 7 PRs.

Include the following sample comprehensive ASCII diagrams verbatim:
\begin{itemize}[leftmargin=*]
\item day\_01/diagrams/agent\_workload\_hour\_000.txt (initialization)
\item day\_03/diagrams/agent\_workload\_hour\_071.txt (safety day end)
\item day\_05/diagrams/agent\_workload\_hour\_119.txt (analysis day end)
\item day\_07/diagrams/agent\_workload\_hour\_167.txt (closeout)
\item one cumulative 7-day text diagram drawn from
sponsor\_168h\_summary.json
\end{itemize}

Identify the following individual day-level and hour-level patient and
robot examples by name:
\begin{itemize}[leftmargin=*]
\item Day 1 (H000-H023, 168 patients, PSL 63.4 to 64.8): name the
initialization patient cohort and the first day's escalations from
day\_01/hourly/output/sponsor\_hour\_000.json through
sponsor\_hour\_023.json.
\item Day 3 (H048-H071, 218 patients, PSL 66.0 to 67.0): name the safety
day events and the robot\_execution\_gateway authorization counts.
\item Day 5 (H096-H119, 195 patients, PSL 68.0 to 68.8): name the
analysis-day biostats agent decisions.
\item Day 7 (H144-H167, 150 patients, PSL 69.5 to 70.0): name the
closeout patient cohort and the audit\_trail\_manager final state.
\end{itemize}

The local verification block (one paragraph, importance critical):
explain that
sponsor/final\_paper/168\_hours/instructions/core\_i5\_6200u\_4gb/README.md
documents a complete 7-day run on a 2015-vintage Intel Core i5-6200U
laptop with 4 GB RAM running Windows 10 Pro, including hardware-specific
mitigations (thermal throttling, Windows Update disable, antivirus
exclusion, pagefile sizing). The fact that a 168-hour autonomous sponsor
simulation completes on this hardware proves that the computational
overhead is not the limiting factor for AI-driven oncology trial patient
prediction, which directly contradicts a common objection that LLM-scale
prediction requires data-center hardware. The sponsor\_168h\_summary.json
output from that run mirrors the cloud run output, and the agent scripts
under day\_NN/hourly/ that were partially verified locally on this
hardware are listed by file name. Tie this to the FDA RTCT vision: a
small hospital site running a Core i5-6200U could in principle stream
sponsor-side decisions to the FDA in real time without provisioning new
infrastructure.

Advantages (one paragraph): Simulation 4 extends Simulation 3 by a
factor of 7, demonstrating multi-day continuity, week-scale PSL
improvement (+6.6), and the largest agent-script footprint of the four
simulations. The 168 GitHub commits across 7 branches and 7 PRs prove
that an autonomous sponsor can also operate the version-control workflow
of a real pharmaceutical organization.

Disadvantages (one paragraph): like Simulation 3, this simulation lacks
the minute-resolution narrative that Simulation 1 provides. The
verification on a 4 GB laptop, while important, was partial - not every
script executed end-to-end on that hardware, and downstream analytics
must account for the partial-verification status.

Cite \cite{kawchak_2026_19396256} and \cite{fda2026realtime}.]

\subsection{Cross-Simulation Synthesis}
\label{sec:results-synthesis}

[PROCESSING INSTRUCTION (Synthesis, 2 to 3 paragraphs and 1 comprehensive
ASCII comparison diagram): Compose a short closing subsection that draws
the cross-simulation comparison. Render a single ASCII comparison diagram
in a verbatim block with four columns (Simulation 1 through 4) and rows
covering: scope, hours, patients, robots, ASCII diagrams generated,
Python scripts generated, JSON outputs generated, GitHub commits, and
local verification status. Source the row values from the per-simulation
subsections above. Close with one sentence stating that the cross-
simulation pattern - repository-scale 1M-token context, hourly commit
cadence, and ASCII-diagram production - is the computational signature
that distinguishes this work from the supervised baselines named in the
Introduction.]
